% Hand-written; appended to tables.tex by scripts/trigger_latex_tables.py.
\clearpage
\section*{Conclusions: resampling triggers}

\paragraph{Setup.} Each trigger was run through the same harness, observation sequences, propagation noise and resampler seeds as the published $E_2<0.5N$ runs, so every trial is compared with its own $E_2$ counterpart. The comparison covers 6 filtering models, $N\in\{32,128,512,2048\}$ ($N\le128$ for OT and Gumbel) and 15 resamplers. The common-threshold and matched-rate sets add 672 cells each, 1344 in total; none failed, and no trial had a numerical failure. A cell counts as \emph{better} or \emph{worse} when the 95\% CI of the paired difference, clustered by observation sequence, excludes 0.

\paragraph{How often each trigger resamples.} At the common threshold $0.5N$ the measures are ordered $E_\infty\le E_2\le E_H$, so they fire at different rates. Relative to $E_2$, $E_H$ resamples 4--29\% less often and $E_\infty$ 1.0--3.0$\times$ as often; the largest change is on stochastic volatility and the smallest on Lorenz-96, which resamples at nearly every step under every trigger. $CV^2>c$ adds no new rule: $CV^2=N/E_2-1$, so $CV^2>1$ is exactly $E_2<0.5N$, and its trials are identical to the baseline.

\paragraph{Standard resamplers are insensitive to the trigger.} The standard methods are systematic, multinomial, stop-gradient, soft, OT, eigen and Gaussian moment matching, Gumbel-softmax and variational KL. For these methods, only 4.4\% of (cell, metric) comparisons are significant. That is the false-positive rate expected at 95\% confidence, and the significant cells are scattered with no consistent direction: each method ties $E_2$ in 20--24 of its 24 cells (11--12 of 12 for OT and Gumbel) on RMSE, predictive NLL and coverage. Within this range of resampling rates, the choice among $E_2$, $E_H$ and $E_\infty$ does not change their accuracy.

\paragraph{DSR is sensitive, and less resampling helps.} For the four DSR variants, 20\% of comparisons are significant, and their direction is consistent.
\begin{itemize}
\item \emph{Resampling more often ($E_\infty$) hurts DSR.} Bearing-only tracking: predictive NLL worsens by 0.49--1.36 nats, significant in 13 of 16 DSR cells, and paper-DSR coverage falls from 0.087 to 0.022 at $N{=}512$. Lorenz-63: coverage drops 3--11 points, significant in 15 of 16 cells. Lorenz-96: RMSE and NLL become slightly worse ($<1\%$).
\item \emph{Resampling less often ($E_H$) helps DSR modestly.} Lorenz-96: RMSE improves by 0.7--4.4\% and NLL by up to 1.07 nats, each significant in 11 of 16 cells. Bearing at $N{=}512$: paper-DSR RMSE improves from 4.87 to 4.67. Bearing NLL also improves for the unsorted and $\tau{=}0.5$ variants.
\item \emph{Rescaled DSR ($\tau{=}2.5/N$) is much less sensitive.} The exception is bearing coverage, where $E_\infty$ still costs 8--12 points at $N\ge128$.
\end{itemize}
This fits the static and gradient results. Each DSR call contracts the particle cloud, so the error is paid per resampling step, and a trigger that fires more often accumulates more of it. For DSR the trigger behaves like a tuning parameter; for unbiased resamplers it does not.

\paragraph{At a matched resampling rate, the measure barely matters.} To separate the measure from the rate, the $E_H$ and $E_\infty$ thresholds were recalibrated per model so that systematic resampling resamples as often as under $E_2<0.5N$ ($\tau_H=0.57$--$0.62$, $\tau_\infty=0.19$--$0.35$). For the standard resamplers, rates then stay within 6\% of $E_2$ in every model except stochastic volatility under $E_H$ (11\% fewer resamplings). They tie $E_2$ in 97--98\% of comparisons, with significant cells split evenly between better and worse, which is a null result. For DSR the share of significant comparisons falls from 18--23\% to 8--9\%, and the effects shrink. On Lorenz-63, the 15 of 16 coverage losses under $E_\infty$ become 3 better and 3 worse. On Lorenz-96, the $E_H$ RMSE gain falls from 0.7--4.4\% to 0.2--1.2\%. The differences that remain follow DSR's own rate, which the systematic-based calibration does not match exactly. On bearing, DSR under matched $E_\infty$ still resamples 17--46\% more often than under $E_2$, and the worst NLL cells (0.34--0.72 nats) are the ones with the largest mismatch. On Lorenz-96, matched $E_H$ resamples 1--2\% less and gives small gains, while matched $E_\infty$ resamples 1--2\% more and gives small losses. We therefore find no evidence that any degeneracy measure is better than another at equal rate. For DSR, the resampling rate is what matters. A side effect is that a threshold calibrated on one resampler does not transfer to DSR: $E_\infty$ fires more often on DSR's contracted clouds than on systematic ones.

\paragraph{What the trigger does not fix.} No trigger closes the gap between paper-DSR and unbiased resampling. On bearing at $N{=}512$, its best coverage (0.113 under $E_H$) remains far from systematic resampling (0.826).

\paragraph{Caveats.} The matched rates are calibrated on systematic resampling with 30 seeds at $N\in\{128,512\}$, so for DSR they are only approximately matched; a calibration on each method's own rate would test the measure more cleanly. On Lorenz-96, coverage is near 0 under every trigger, so the ``better'' coverage verdicts there are changes of about 0.01. The results cover 50-step filters on these six models only. E-MIM is not included, because its exact definition was not available.
